% Original Prism supporting-material draft; no experimental content is supplied.
% Build separately with --main supplementary.tex --no-build-script.
\documentclass[a4paper,11pt]{article}
\usepackage[margin=25mm]{geometry}
\usepackage{amsmath,amssymb,graphicx,array,booktabs,tabularx,siunitx,longtable,newfloat}
\usepackage[version=4]{mhchem}
\usepackage[numbers,sort&compress,square]{natbib}
\usepackage[hidelinks]{hyperref}
\providecommand{\doi}[1]{\href{https://doi.org/#1}{\nolinkurl{#1}}}
\DeclareFloatingEnvironment[fileext=los,name=Scheme,placement=htbp]{scheme}
\sisetup{uncertainty-mode=separate}
\setlength{\emergencystretch}{2em}
\renewcommand{\thepage}{S\arabic{page}}
\renewcommand{\thefigure}{S\arabic{figure}}
\renewcommand{\thetable}{S\arabic{table}}
\renewcommand{\theequation}{S\arabic{equation}}
\renewcommand{\thescheme}{S\arabic{scheme}}
\raggedbottom

\title{Supporting Information: Manuscript title}
\author{First Author and Second Author\\
  \small Department, Institution, City, Country\\
  \small Correspondence: \texttt{second.author@example.org}}
\date{}

\begin{document}
\maketitle
This independent file supports either a research manuscript or a review.
Replace its examples with actual supporting material and remove unused
sections. No experiment, literature search or scientific result is supplied
by these resources. All sample values are invented solely for formatting.
\tableofcontents
\clearpage

\section{Supporting overview}\label{sec:si-overview}
Describe the contents and their relation to the main text. For a research
paper, include only actual supporting evidence and information appropriate
to this file. For a review, use verified evidence tables or additional
analysis only when they were actually prepared. The public reference is a
bibliography demonstration, not evidence for a dye study~\cite{peng2011reproducible}.

Figure~\ref{fig:si-overview}, Scheme~\ref{sch:si-sequence} and
Equations~\eqref{eq:si-ratio}--\eqref{eq:si-difference} demonstrate independent
S numbering. Define actual quantities before using the equations.
\begin{align}
 q_i &= \frac{x_i}{x_{\mathrm{ref}}},\label{eq:si-ratio}\\
 \Delta q &= q_{\mathrm{A}}-q_{\mathrm{B}}.\label{eq:si-difference}
\end{align}
These expressions are layout examples, not a validated optical quantity or method.
\begin{figure}[htbp]
\centering
\fbox{\parbox[c][30mm][c]{0.82\linewidth}{\centering
Original supporting graphic slot\\[7pt]
Replace with the authors' verified figure and legend.}}
\caption{\textbf{Supporting layout example.} No data are shown. Describe the
replacement figure, symbols, units and uncertainty accurately. Preserve
compound and panel identifiers from the main manuscript.}
\label{fig:si-overview}
\end{figure}
\begin{scheme}[htbp]
\centering
\fbox{\parbox[c][18mm][c]{0.82\linewidth}{\centering
\ce{A ->[replace] B}\\[5pt]Synthetic letter slots, not a reported reaction.}}
\caption{A synthetic scheme placeholder; supply original structures and
verified information, or remove it. Formulae such as \ce{CO2} illustrate
chemical typography only.}
\label{sch:si-sequence}
\end{scheme}
\clearpage

\section{Supporting tables}\label{sec:si-tables}
Table~\ref{tab:si-values} demonstrates aligned numbers and uncertainty.
Preserve precision and units when replacing it with actual supporting data.
\begin{table}[htbp]
\centering\small
\caption{Invented values illustrating numerical alignment only.}
\label{tab:si-values}
\begin{tabularx}{\linewidth}{@{}c >{\raggedright\arraybackslash}X
 S[table-format=3.1(2)] S[table-format=1.2(2)]@{}}
\toprule
Compound & Substituent slot & {\shortstack{Illustrative wavelength\\/ nm}}
 & {\shortstack{Illustrative value}}\\
\midrule
\textbf{1} & \ce{CH3} & 420.0 +- 0.6 & 0.31 +- 0.04\\
\textbf{2} & \ce{C2H5} & 505.0 +- 0.4 & 0.28 +- 0.02\\
\textbf{3} & \ce{F} & 478.0 +- 0.8 & 0.19 +- 0.09\\
\bottomrule
\end{tabularx}
\par\smallskip\footnotesize\raggedright
All values are invented. Define the actual endpoint, units, uncertainty
measure and replicate counts; no measurement was performed.
\end{table}

Table~\ref{tab:si-inventory} demonstrates a long table whose column headings
repeat if it spans pages. Its entries are writing prompts, not an inventory of completed
experiments or a review's search process. Replace or delete it as needed.
\begingroup\small
\begin{longtable}{@{}p{0.12\linewidth}p{0.24\linewidth}>{\raggedright\arraybackslash}p{0.54\linewidth}@{}}
\caption{Synthetic supporting inventory used to demonstrate continuation.}
\label{tab:si-inventory}\\
\toprule Item & Topic & Description\\\midrule
\endfirsthead
\multicolumn{3}{@{}l}{Table~\thetable\ (continued)}\\[3pt]
\toprule Item & Topic & Description\\\midrule
\endhead
\midrule\multicolumn{3}{r@{}}{Continued on next page}\\
\endfoot
\bottomrule
\endlastfoot
S01 & Scope & Explain which actual material supports the main manuscript.\\
S02 & Identifiers & Preserve the compound identifiers used in the main text.\\
S03 & Provenance & Attribute actual data or cited material to its source.\\
S04 & Structure & Supply verified structural information when applicable.\\
S05 & Characterization & Link actual characterization evidence to its compound.\\
S06 & Conditions & State relevant conditions without inventing experiments.\\
S07 & Controls & Describe controls actually used, where applicable.\\
S08 & Measurement & Define the endpoint represented by each reported value.\\
S09 & Units & Preserve original units and explain any actual conversions.\\
S10 & Precision & Preserve the numerical precision justified by the evidence.\\
S11 & Uncertainty & Define the meaning and source of any uncertainty values.\\
S12 & Replicates & State the actual number and nature of replicates.\\
S13 & Missing values & Distinguish unavailable data from a measured zero.\\
S14 & Software & Record relevant versions and settings for reproducibility.\\
S15 & Comparability & Explain differences that limit a comparison.\\
S16 & References & Verify the identity of cited publications and their links.\\
S17 & Figures & Keep original editable graphics and readable final-size labels.\\
S18 & Schemes & Preserve structures, identifiers and verified conditions.\\
S19 & Tables & Relate supporting entries to the corresponding main-text claim.\\
S20 & Permissions & Record permissions where material is reused.\\
S21 & Evidence & Separate observations from interpretation or proposals.\\
S22 & Limitations & Report actual limitations and unresolved questions.\\
S23 & Availability & Describe real access conditions for underlying material.\\
S24 & Consistency & Check labels and terminology against the main manuscript.\\
\end{longtable}
\endgroup

\section{Additional information}\label{sec:si-additional}
Use this section only for supporting information actually prepared by the
authors. It is not a prefilled experimental protocol or literature-search
record. Remove unused examples and instructional text before submission.

\bibliographystyle{elsarticle-num}
\bibliography{references}
\end{document}
